\section{Review of Previous Work}
\label{sec:literature_review}

The integration of blockchain technology into UAV communication networks has been an active area of research, with numerous studies exploring its potential to enhance security, trust, and decentralization. This section provides a critical review of the state-of-the-art, focusing on key research pillars including routing, authentication, consensus mechanisms, and overall security architectures.

A comprehensive survey by Hafeez et al. \cite{hafeez2023blockchain} provides a foundational map of the landscape, categorizing blockchain use cases in UAV communications and highlighting the trade-offs between different blockchain models. The survey underscores the nascent stage of the field and calls for more research into lightweight consensus protocols and cross-layer optimization, a call that motivates the work presented in this paper.

\subsection{Routing and Trust Management}
Secure routing is paramount in FANETs. Several recent works have leveraged blockchain to establish trust and ensure path integrity. Jia et al. \cite{jia2025trusted} propose a trusted routing scheme using multi-agent deep reinforcement learning (MARL), where blockchain is employed to manage a trust evaluation system for UAV nodes. Similarly, He et al. \cite{he2025trusted} introduce a zero-trust architecture for Low-Altitude Internet of Networks (LAINs), utilizing a permissioned blockchain for routing optimization and trust establishment. A notable contribution by Luo et al. \cite{luo2023escm} is the ESCM framework, which co-designs a routing protocol based on the artificial bee colony (ABC) algorithm with a novel Proof-of-Network-Coding (PoNC) blockchain consensus mechanism. This integrated approach aims to accelerate communication while securing the network, although the overhead of PoNC in highly dense networks remains an area for further investigation.

\subsection{Authentication and Identity Management}
Securely authenticating UAVs is a critical prerequisite for any secure communication protocol. Traditional PKI systems often rely on centralized certificate authorities, which are unsuitable for decentralized FANETs. To address this, recent research has focused on innovative, lightweight authentication schemes. Jiao et al. \cite{jiao2025lightweight} introduce a novel scheme that integrates blockchain with aggregatable subvector commitments (aSVC) to compress identity states. This method significantly reduces on-chain storage to a constant level and optimizes the computational complexity of batch authentication, making it highly suitable for resource-constrained UAV swarms. In a different approach, Hafeez et al. \cite{hafeez2024beta} developed BETA-UAV, a blockchain-based authentication scheme using Elliptic Curve Cryptography (ECC) and smart contracts for logging authentication events. Further extending the concept of authentication to physical payloads, Ajakwe et al. \cite{ajakwe2025ibanda} propose the iBANDA framework, which combines AI-based visual detection (YOLOv5s) with a Snowball PoS consensus mechanism for dual-layer authentication of both the drone and its cargo in logistics applications. This hybrid approach demonstrates the potential of fusing AI with blockchain to create more robust and context-aware security solutions.

\subsection{Consensus Mechanisms for FANETs}
The choice of consensus mechanism is critical in determining the performance and scalability of a blockchain-based UAV network. Hossain et al. \cite{hossain2024blockchain} provide a valuable performance analysis of private blockchain implementations, evaluating metrics such as throughput and latency for PBFT-like algorithms. Their findings confirm that while private blockchains can offer low latency, their performance degrades significantly as the number of nodes increases. This highlights the need for more scalable consensus protocols. Addressing this, Huang et al. \cite{huang2025performance} analyze the performance of chained-HotStuff, a modern BFT consensus protocol, under realistic UAV network conditions, providing insights into its viability. The development of truly lightweight and energy-efficient consensus mechanisms, such as the PoC protocol in our PrimeFusion framework, remains a significant and active research challenge.

\subsection{Comprehensive Security Architectures and Applications}
Several researchers have proposed holistic security architectures that integrate blockchain across multiple layers of the communication stack. Jagatheesaperumal et al. \cite{jagatheesaperumal2023blockchain} present a comprehensive security architecture for UAVs in B5G/6G networks, where blockchain serves as a foundational layer for security services. Their work emphasizes a cross-layer approach, integrating blockchain with physical layer security and AI-driven threat detection. In a specific application context, Khor et al. \cite{khor2025secure} propose a secure communication protocol for public Unmanned Traffic Management (UTM) using LoRa for long-range, low-power communication, with a public blockchain for collision avoidance and traffic management. While these comprehensive architectures are powerful, they often come with significant implementation complexity. The PrimeFusion framework, presented in the following sections, aims to address this gap by proposing a more lightweight and modular architecture that balances security with performance.

